\ResearchStyle
\PartDivider{PART II}{Research Problem and Gap}{Reconstruct the exact unfairness problem that motivates CATRA.}

\begin{frame}{Paper Context and Scope}
\PaperClaim{The paper studies contention between stations and between flows in IEEE 802.11 multi-hop ad hoc networks.}
\begin{columns}[T,onlytextwidth]
\column{0.48\textwidth}
\textbf{Included}
\begin{itemize}
\item Shared IEEE 802.11 medium.
\item Multi-hop forwarding.
\item TCP flow contention.
\item MAC and transport feedback.
\end{itemize}
\column{0.48\textwidth}
\textbf{Explicitly excluded}
\begin{itemize}
\item Mobility-induced unfairness.
\item Multi-rate operation.
\item Multi-channel operation.
\end{itemize}
\end{columns}
\end{frame}

\begin{frame}{Problem 1: MAC-Layer Unfairness}
\PaperClaim{BEB may be unfair in multi-hop ad hoc networks: CW is chosen mainly from congestion conditions.}
\begin{center}
\begin{tikzpicture}[node distance=0.42cm,
flowbox/.style={draw=ResearchColor,rounded corners,fill=CardBg,text width=2.65cm,minimum height=0.9cm,align=center}]
\node[flowbox] (beb) {IEEE 802.11 BEB};
\node[flowbox,right=of beb] (state) {CW follows congestion conditions};
\node[flowbox,right=of state] (missing) {Neighbor and higher-layer conditions are not considered};
\node[flowbox,draw=WarningColor,right=of missing] (unfair) {CW may be unsuitable for fairness};
\draw[-{Latex[length=2mm]}] (beb)--(state);
\draw[-{Latex[length=2mm]}] (state)--(missing);
\draw[-{Latex[length=2mm]},draw=WarningColor] (missing)--(unfair);
\end{tikzpicture}
\end{center}
\vspace{0.15cm}
\begin{itemize}
\item BEB appears to give contending stations the same channel-access opportunity.
\item It does not account for the number of flows in the channel or users in the system.
\item Consequently, disadvantaged stations may have difficulty accessing the channel.
\end{itemize}
\end{frame}

\begin{frame}{Problem 2: TCP-Layer Unfairness}
\PaperClaim{An excessively large TCP congestion window creates too many packets competing for the shared medium.}
\begin{center}
\begin{tikzpicture}[node distance=0.48cm,
flowbox/.style={draw=ResearchColor,rounded corners,fill=CardBg,text width=2.6cm,minimum height=0.85cm,align=center}]
\node[flowbox] (grow) {TCP AIMD grows $cwnd$};
\node[flowbox,right=of grow] (packets) {Too many packets compete};
\node[flowbox,draw=WarningColor,right=of packets] (cont) {Network congestion increases};
\node[flowbox,draw=WarningColor,right=of cont] (fair) {Throughput and fairness degrade};
\draw[-{Latex[length=2mm]}] (grow)--(packets);
\draw[-{Latex[length=2mm]},draw=WarningColor] (packets)--(cont);
\draw[-{Latex[length=2mm]}] (cont)--(fair);
\end{tikzpicture}
\end{center}
\vspace{0.18cm}
\begin{itemize}
\item TCP performance in a multi-hop wireless network depends critically on $cwnd$.
\item The paper observes that some flows receive very large bandwidth while others receive little bandwidth.
\item Conventional TCP rate control can therefore produce unfair bandwidth sharing among flows.
\end{itemize}
\end{frame}

\begin{frame}{Problem Statement}
\begin{center}
\textbf{TCP traffic unfairness is due to both DCF channel access and TCP traffic-rate control.}
\end{center}
\begin{columns}[T,onlytextwidth]
\column{0.48\textwidth}
\textbf{DCF channel-access behavior}
\begin{itemize}
    \item Some stations may have difficulty accessing the channel.
\end{itemize}

\column{0.48\textwidth}
\textbf{TCP traffic-rate behavior}
\begin{itemize}
    \item Some TCP flows obtain very large bandwidth, while other flows obtain little bandwidth.
\end{itemize}
\end{columns}

\begin{center}
\begin{tikzpicture}[node distance=0.55cm]
\node[draw=ResearchColor,rounded corners,fill=CardBg,text width=3.1cm,align=center] (mac) {DCF channel-access problem};
\node[font=\bfseries,right=of mac] (plus) {$+$};
\node[draw=ResearchColor,rounded corners,fill=CardBg,text width=3.1cm,align=center,right=of plus] (tcp) {TCP traffic-rate problem};
\node[draw=WarningColor,rounded corners,fill=CardBg,text width=3.8cm,align=center,right=of tcp] (fair) {Unfair TCP traffic};
\draw[-{Latex[length=2.5mm]},draw=WarningColor] (tcp)--(fair);
\end{tikzpicture}
\end{center}
\end{frame}

\begin{frame}{Gap Identified from Related Work}
\begin{center}
\textbf{The reviewed approaches leave limitations when station and flow fairness are considered together.}
\end{center}

\begin{columns}[T,onlytextwidth]
\column{0.32\textwidth}
\textbf{MAC approaches}
\begin{itemize}
\item EIFS, back-off, and CW adaptation.
\item Some need topology-dependent values; most per-flow CW studies are single-hop.
\end{itemize}

\column{0.32\textwidth}
\textbf{TCP approaches}
\begin{itemize}
\item NRED and TCP window or rate adjustment.
\item Primarily control TCP behavior.
\end{itemize}

\column{0.32\textwidth}
\textbf{Cross-layer approaches}
\begin{itemize}
\item Use MAC information to adapt TCP.
\item A cited method omits MAC contention and cannot exactly identify intermediate-node efficiency.
\end{itemize}
\end{columns}

\vspace{-0.05cm}
\begin{center}
\begin{tikzpicture}[node distance=0.55cm]
\node[draw=MethodColor,rounded corners,fill=CardBg,text width=4.2cm,align=center,inner ysep=3pt] (cw) {Suitable CW for each station};
\node[draw=MethodColor,rounded corners,fill=CardBg,text width=4.2cm,align=center,inner ysep=3pt,right=of cw] (rate) {TCP sending rate for each flow};
\node[draw=WarningColor,rounded corners,fill=CardBg,right=of rate] (catra) {CATRA};
\draw[-{Latex[length=2.2mm]},draw=MethodColor] (cw)--(rate);
\draw[-{Latex[length=2.2mm]},draw=MethodColor] (rate)--(catra);
\end{tikzpicture}
\end{center}
\end{frame}


\begin{frame}{Research Need: Two-Stage Control}
\begin{center}
\begin{tikzpicture}[node distance=0.8cm]
\node[draw=ResearchColor,rounded corners,fill=CardBg,text width=4.2cm,align=center,minimum height=1.1cm] (s1)
{Stage 1\\Determine a suitable CW for each station};
\node[draw=ResearchColor,rounded corners,fill=CardBg,text width=4.2cm,align=center,minimum height=1.1cm,right=of s1] (s2)
{Stage 2\\Regulate the sending rate of each TCP flow};
\draw[-{Latex[length=2.8mm]},draw=ResearchColor,line width=1.2pt] (s1)--(s2);
\end{tikzpicture}
\end{center}
\vspace{0.5cm}
\begin{itemize}
\item Stage 1 targets fair channel access between stations.
\item Stage 2 targets fair bandwidth sharing between end-to-end flows.
\item Both stages use channel-utilization information collected at MAC/PHY.
\end{itemize}
\end{frame}

\input{slides/13-research-questions}

\begin{frame}{Main Contributions of the Paper}
\begin{enumerate}
\item Define station-level and flow-level Fair Bandwidth Ratios.
\item Estimate real channel utilization from packet airtime.
\item Adapt the IEEE 802.11 contention window using $RBR_S/FBR_S$.
\item Pass MAC/PHY utilization information to the transport layer.
\item Adapt TCP generation delay and $cwnd$ using $RBR_f/FBR_f$.
\end{enumerate}
\vspace{0.35cm}
\KeyPoint{CATRA is a feedback-control design with separate actuators for station contention and flow demand.}
\end{frame}
